\documentclass{article}
\usepackage[UTF8]{ctex}
\usepackage{amsmath,amssymb,amsfonts}
\usepackage{enumitem}
\usepackage{geometry}
\geometry{a4paper, margin=1in}
% === 错题本排版优化（不参与题目标识） ===
\setlength{\parindent}{0pt}
\setlength{\parskip}{0.35em plus 0.10em minus 0.08em}
\setlength{\emergencystretch}{3em}
\setlist{itemsep=0.25em, topsep=0.40em, parsep=0pt}
\allowdisplaybreaks[2]
% === 错题本排版优化结束 ===
\begin{document}
\section*{陕西师范大学2024年高等代数真题}
\section*{177．陕西师范大学2024年高等代数试卷}

\subsection{题目1}
一．（15 分）已知 \(n(n \geq 2)\) 阶行列式

\[
D=\left|\begin{array}{cccccc}
1 & 2 & 3 & \cdots & n-1 & n \\
1 & -1 & 0 & \cdots & 0 & 0 \\
0 & 2 & -2 & \cdots & 0 & 0 \\
\vdots & \vdots & \vdots & & \vdots & \vdots \\
0 & 0 & 0 & \cdots & 2-n & 0 \\
0 & 0 & 0 & \cdots & n-1 & 1-n
\end{array}\right| .
\]

\(A_{i j}\) 为 \(D\) 的第 \(i\) 行 \(j\) 列元素 \(a_{i j}\) 的代数余子式，求 \(A_{11}+A_{12}+\cdots+A_{1 n}\) 。

\subsection{题目2}
二．（15 分）证明复数域上的多项式

\[
f(x)=x^{n}+n x^{n-1}+(n(n-1)) x^{n-2}+\cdots+(n(n-1)(n-2) \cdots 4 \cdot 3) x^{2}+(n!) x+n!
\]

没有重根。

\subsection{题目3}
三．（20 分）设 \(\alpha_{1}, \alpha_{2}, \cdots, \alpha_{m}\) 是线性方程组 \(A X=0\) 的一个基础解系，记

\[
\beta_{1}=t_{1} \alpha_{1}+t_{2} \alpha_{2}, \beta_{2}=t_{1} \alpha_{2}+t_{2} \alpha_{3}, \cdots, \beta_{m}=t_{1} \alpha_{m}+t_{2} \alpha_{1} .
\]

其中 \(t_{1}, t_{2}\) 为常数，求 \(t_{1}, t_{2}\) 满足何种关系时，\(\beta_{1}, \beta_{2}, \cdots, \beta_{m}\) 也为方程组 \(A X=0\) 的一个基础解系。

\subsection{题目4}
四．（15 分）设 \(A\) 是秩为 \(r\) 的 \(n\) 阶方阵，证明：存在 \(n\) 阶可逆矩阵 \(P\) ，使得 \(P A P^{-1}\) 的后 \(n-r\) 行全为零。

\subsection{题目5}
五．（20分）设 \(A=\left(a_{i j}\right)\) 是 \(m \times n\) 阶实矩阵，证明：\(n\) 元实二次型

\[
f\left(x_{1}, x_{2}, \cdots, x_{n}\right)=\left(a_{11} x_{1}+a_{12} x_{2}+\cdots+a_{1 n} x_{n}\right)^{2}+\cdots+\left(a_{m 1} x_{1}+a_{m 2} x_{2}+\cdots+a_{m n} x_{n}\right)^{2}
\]

是正定二次型的充要必要条件是 \(A\) 的秩等于 \(n\) 。

\subsection{题目6}
六．（15 分）设 \(A=\binom{A_{1}}{A_{2}}\) 是数域 \(P\) 上的 \(n(n \geq 2)\) 阶可逆分块矩阵，记

\[
W_{1}=\left\{X \in P^{n} \mid A_{1} X=0\right\}, W_{2}=\left\{X \in P^{n} \mid A_{2} X=0\right\} .
\]

证明：\(P^{n}=W_{1} \oplus W_{2}\) ．

\subsection{题目7}
七．（20 分）设 \(\sigma\) 和 \(\tau\) 是数域 \(P\) 上 \(n\) 维线性空间 \(V\) 上的两个线性变换，\(\sigma\) 在数域 \(P\) 上有 \(n\) 个互不相同的特征值，证明：\(\sigma\) 的特征向量都是 \(\tau\) 的特征向量的充要必要条件是 \(\sigma \tau=\tau \sigma\) ．

\subsection{题目8}
八．（15 分）已知复数域上的三阶矩阵

\[
A=\left(\begin{array}{lll}
2 & 0 & 0 \\
a & 1 & 0 \\
b & c & 1
\end{array}\right)
\]

求 \(A\) 所有可能的 Jordan 标准形，并确定 \(A\) 可对角化的条件．

\subsection{题目9}
九．（15 分）设 \(A\) 与 \(B\) 都是 \(n\) 阶正定矩阵，证明：\(A B\) 是正定矩阵的充要条件是 \(A B=B A\) ．
\end{document}
